% Generated draft: requires mathematical and source review.
\documentclass[11pt]{article}
\usepackage[margin=25mm]{geometry}
\usepackage{fontspec}
\setmainfont{Latin Modern Roman}
\usepackage{amsmath,amssymb,mathtools,unicode-math}
\setmathfont{Latin Modern Math}
\usepackage{enumitem,needspace,xurl,hyperref}
\setlength{\parindent}{0pt}
\setlength{\parskip}{5pt}
\setlength{\emergencystretch}{3em}
\newcommand{\N}{\mathbb{N}}
\newcommand{\Z}{\mathbb{Z}}
\newcommand{\Q}{\mathbb{Q}}
\newcommand{\R}{\mathbb{R}}
\newcommand{\C}{\mathbb{C}}
\newcommand{\F}{\mathbb{F}}
\newcommand{\A}{\mathbb{A}}
\newcommand{\PP}{\mathbb{P}}
\newcommand{\E}{\mathbb{E}}
\newcommand{\eps}{\varepsilon}
\newcommand{\dd}{\,\mathrm{d}}
\begin{document}
\section*{UnsolvedMath ID: 800414}
\textbf{Source problem code:} catalogue entry
\subsection{Definitions and mathematical statement}
Use the standard non-Archimedean interpretation of the source notation: \(\mathbb C_v\) is a complete algebraically closed non-Archimedean valued field, and \(\mathsf P^1_{\mathrm{Berk}}\) is its Berkovich projective line. This ambient-field interpretation is implicit in the notation rather than explicitly stated in the supplied extract. Let \(K\subset\mathbb C_v\) be a subfield on which the restricted valuation is discrete, and suppose that an algebraic closure \(\bar K\subseteq\mathbb C_v\) is dense in \(\mathbb C_v\). Let \(p\geq 0\) be the characteristic of the residue field of \(\mathbb C_v\), and let \(f\in\mathbb C_v(z)\) have degree at least \(2\). In particular, the supplied statement does not require \(f\in K(z)\).

For \(\zeta\in\mathsf P^1_{\mathrm{Berk}}\), write \(\deg_\zeta(f)\) for the positive integer local degree of \(f\) at \(\zeta\). The local-degree hypothesis is
\[
 p=0\quad\text{or}\quad
 \bigl(p>0\ \text{and}\ \forall\zeta\in\mathsf P^1_{\mathrm{Berk}},\ p\nmid\deg_\zeta(f)\bigr).
\]
Here the residue-characteristic-zero case imposes no local-degree restriction; this makes explicit the convention implicit in the source's divisibility wording.

Write \(J(f)\) for the Julia set and \(F(f)\) for its complement. Under the Berkovich-component interpretation, a wandering domain is a connected component \(U\) of \(F(f)\) such that the components \(U_n\) containing \(f^n(U)\), for \(n\geq0\), are pairwise distinct. A Type II point is a Berkovich point represented by a closed disc of radius in \(\lvert\mathbb C_v^\times\rvert\). A periodic point \(\zeta\) satisfies \(f^q(\zeta)=\zeta\) for some integer \(q\geq1\).

The supplied Conjecture 37 asserts that, under the hypotheses above, every wandering domain \(U\) satisfies
\[
 U\subseteq\mathcal B_f(\zeta)
 \quad\text{for some periodic Type II point }\zeta\in J(f),
\]
where \(\mathcal B_f(\zeta)\) denotes the source's attracting basin of that periodic point. The supplied extract does not specify the topology or convergence convention defining this basin, or explicitly identify which notion of Fatou component is intended. Thus the component interpretation above is provisional, and the basin terminology is retained without imposing an additional definition. In particular, no multiplier condition or usual attracting-cycle classification is added to the stated Julia-point conclusion.
\subsection{Short English statement}
The supplied conjecture says that, when no Berkovich local degree is divisible by the positive residue characteristic, every wandering domain lies in the attracting basin of a periodic Type II Julia point. In residue characteristic zero the local-degree condition is vacuous. The precise domain and basin conventions require review against the full source.
\subsection{Sources}
\begin{itemize}
\item Local UnsolvedMath registry, permanent ID \(800414\), source problem code \texttt{catalogue entry}, title \emph{Open Problems in Arithmetic Dynamics --- Conjecture 37}. The supplied registry labels this record \emph{open}; that label is source metadata, not an independently verified status assessment.
\item Benedetto, DeMarco, Ingram, Jones et al., \emph{Current Trends and Open Problems in Arithmetic Dynamics} (2018), Conjecture 37. Supplied source URL: \url{https://arxiv.org/abs/1806.04980}. Attribution and proposed year are reproduced from the supplied record.
\item The supplied provenance reports access on 2026-10-16 and extraction method \texttt{arxiv\_latex}, with upstream revision \texttt{d371ee917ac9647fbb06cd6fc529deec6f3dda82}. These are supplied provenance claims; no source visit or independent verification is asserted here.
\end{itemize}
\end{document}
